\section{第五日第一讲}

\subsection{回顾}

回忆我们想要在种种设定（整体算术、整体几何、局部算术、局部几何）下等同两个理论，即（$G$ 的）自守理论与（$\check G$ 的）谱理论。特别地，边界理论（例如 L-函数与周期）应当匹配。举例来说，Fargues--Scholze 关于算术局部域（如 $\mathbb Q_p$）上局部 Langlands 对应的表述，与曲线上几何 Langlands 有相似之处。

在局部算术设定中，回忆 $C_c^\infty(X_G)^{G_O}$ 对应对偶空间 $\check M$ 及其上的 $\check G$ 作用。把它范畴化，就得到局部几何猜想，它联系
\[
  \Shv(X_F/G_O)\cong\QCoh(\check M/\check G).
\]

\begin{exmp}
对群情形，即 $G\times G$ 作用在 $X=G$ 上，在对偶侧有 $\check G\times\check G$ 作用在 $\check X=\check G$ 上，于是 $\check M=T^*\check G\cong\check G\times\check{\mathfrak g}^*$。局部几何猜想此时还原为导出几何 Satake 等价：A 侧是 $\Shv(G_O\backslash G_F/G_O)$，B 侧是 $\QCoh(\check{\mathfrak g}^*/\check G)$。B 侧经 Koszul 对偶与 TQFT 在 raviolo 上的赋值等同。更精确地说，我们有
\[
  B_{\check G}\Big(D\coprod_{D^\times}D\Big)=\QCoh^!(\check{\mathfrak g}[-1]/\check G)\cong\QCoh(\check{\mathfrak g}^*/\check G).
\]
这里第一个恒等式来自 TQFT 的场描述，最后一个等价是 Koszul 对偶。
\end{exmp}

今天的目标是讨论整体几何猜想，然后理解把整体几何猜想与局部几何猜想、整体算术猜想联系起来的那些箭头。

\subsection{设定}

固定 $\Sigma$ 为 $\mathbb C$ 或 $\mathbb F_q$ 上光滑射影代数曲线。

有群 $G$ 作用在 $X$ 上，我们给它配上周期层 $P_X\in\Shv(\Bun_G\Sigma)$。对偶地，$\check G$ 作用在 $\check X$ 上，我们给它配上 L-层 $L_{\check X}\in\QCoh^!(\Loc_{\check G}\Sigma)$。无分歧几何 Langlands 匹配这两个范畴，而相对 Langlands 进一步匹配对象，也就是说，周期层应当匹配 L-层。

\begin{rem}
相对朗兰兹对偶给出的对象匹配，比几何 Langlands 对应的表述与证明中所用的多得多。
\end{rem}

\subsection{$\Bun^X_G$ 与周期层}

考虑 $\Bun^X_G=\Bun^X_G(\Sigma)=\Maps(\Sigma,X/G)$，它是分类「$\Sigma$ 上的 $G$-丛连同相伴 $X$-丛的截面」的栈。这个栈带有一个自然的映射
\setcounter{equation}{2}
\begin{equation}
  \pi:\ \Bun^X_G\to\Bun_G.
  \label{eq:13-3}
\end{equation}
% [Tutor Note: 原文的 (13.1)(13.2) 编号在印刷本中未出现（节内计数器却已进两位），译本按印刷面保留无编号显示，并把计数器置为 2，使 (13.3) 之后的编号与原文一致。]

\begin{defn}
我们由公式
\[
  P_X:=\pi_!\,k_{\Bun^X_G}
\]
定义周期层。
\end{defn}

\begin{exmp}
取 $G=\mathbb G_m$，则
\[
  \Bun_G=\Pic_\Sigma\cong\mathbb Z\times\Jac_\Sigma\times B\mathbb G_m.
\]
\end{exmp}

\begin{exmp}
考虑 $\Bun_{\mathrm{SL}_2}(\mathbb P^1)$。它不紧。它是连通的，其点对应 $\mathcal O\oplus\mathcal O$、$\mathcal O(1)\oplus\mathcal O(-1)$、……。它们被赋予这样的拓扑：当 $0\le b\le a$ 时，$\mathcal O(a)\oplus\mathcal O(-a)$ 落在 $\mathcal O(b)\oplus\mathcal O(-b)$ 的闭包中。
\end{exmp}

\subsection{局部系}

与 $\Bun_G$ 对偶地，我们考虑
\begin{equation}
  \Loc^{\check X}_{\check G}\Sigma:=\Maps_{k\text{-}l.c.}(\Sigma,\check X/\check G),
  \label{eq:13-7}
\end{equation}
它带有一个映射
\begin{equation}
  \check\pi:\ \Loc^{\check X}_{\check G}\Sigma\to\Loc_{\check G}\Sigma:=\Maps_{k\text{-}l.c.}(\Sigma,\mathrm{pt}/\check G),
  \label{eq:13-8}
\end{equation}
其中下标 $k\text{-}l.c.$ 表示局部常值映射。

\begin{exmp}
对 $G=\mathbb G_m$，有
\begin{equation}
  \Loc_{\check G}=B\mathbb G_m\times(\mathbb G_m)^{2g}\times(\text{导出部分}),
  \label{eq:13-10}
\end{equation}
其中 $g$ 是 $\Sigma$ 的亏格。更典范地说，因子 $(\mathbb G_m)^{2g}$ 可以看作 $H^1(\Sigma,\mathbb C^\times)$。
\end{exmp}

\subsection{L-层及其纤维}

我们对 $\check G$-空间 $\check X$ 定义相应的 L-层，并研究它的纤维。回忆我们有映射 $\check\pi:\Loc^{\check X}_{\check G}\to\Loc_{\check G}$。

\begin{defn}
我们由公式
\[
  L_{\check X}:=\check\pi_*\,\omega_{\Loc^{\check X}_{\check G}}\in\QC^!(\Loc_{\check G})
\]
定义 L-层。
\end{defn}

\begin{rem}
层 $\omega_{\Loc^{\check X}_{\check G}}$ 与 $\mathcal O_{\Loc^{\check X}_{\check G}}$ 之间有同构。
\end{rem}

\begin{rem}
为什么要用对偶化层而不是结构层？大致是因为范畴 $\IndCoh$ 偏爱对偶化层，而范畴 $\QCoh$ 偏爱结构层。
\end{rem}

\begin{rem}
我们得对 L-层作归一化，它在函数层面相当于对 L-函数作归一化，以使内积有意义。
\end{rem}

给定局部系 $E\in\Loc_{\check G}$，$\check\pi$ 在 $E$ 处的纤维有如下几种等价描述：

\begin{enumerate}[label=(\arabic*),leftmargin=1.8em]
  \item $E$ 的相伴 $\check X$-丛的局部常值截面；
  \item 被 $E$ 扭的局部常值映射 $\Sigma\to\check X$；
  \item $\check X$ 上 $\pi_1(\Sigma)$ 的不动点。
\end{enumerate}

在最后一种描述里，我们用的是 $\pi_1(\Sigma)$ 经由 $E:\pi_1(\Sigma)\to\check G$ 拉回 $\check G$ 在 $\check X$ 上的作用而得到的作用。事实上，上述所有描述都应在导出意义下理解。把 $\check\pi$ 在 $E$ 处的纤维记作 $\check\pi^{-1}E$。

\begin{rem}
人们会问：给定群 $H$ 作用在空间 $Y$ 上，它的导出不动点轨迹是什么意思？经典地，不动点轨迹可以解释为从 $\mathrm{pt}/H$ 到 $Y$ 的映射空间。我们通过考虑从 $\mathrm{pt}/H$ 到 $Y$ 的导出映射空间来理解导出不动点轨迹。
\end{rem}

由定义，给定局部系 $E$，$L_{\check X}$ 在 $E$ 处的纤维是 $R\Gamma(\check\pi^{-1}E,\omega)$。因此 $L_{\check X}$ 在「$R\Gamma(\check\pi^{-1}E,\omega)$ 有限维」的轨迹上是凝聚的；等价地，在（经典）不动点个数有限处凝聚。这也可以说成（经典）不动点孤立。

\begin{exmp}
设 $\check X$ 是 $\check G$ 的线性表示。不动点孤立等价于 $\pi_1\Sigma$ 在 $\check X$ 上没有平凡子表示。这个条件决定了 $\Loc_{\check G}$ 上使 L-层 $L_{\check X}$ 凝聚的轨迹。
\end{exmp}

\subsection{整体几何猜想}

现在我们陈述 BZSV 的整体几何猜想。先陈述一个朴素版本。

\begin{conj}[朴素整体几何猜想]
周期层 $P_X$ 与 L-层 $L_{\check X}$ 在几何 Langlands 对应下匹配：
\[
  \mathbb L_G:\ \Shv(\Bun_G)\to\QC^!(\Loc_{\check G}),\qquad P_X\mapsto L_{\check X}.
\]
\end{conj}

回忆 A 侧有三种层论可供选择，即 de Rham、Betti 与 étale。除 de Rham 设定外，这个猜想应当作修改：对几何 Langlands 对应的陈述而言，A 侧的层应当满足某个条件 $N$。从 TQFT 的角度看这是自然的。由 TQFT 的局部常值公理，Hecke 范畴（即 $A_G(D_x\coprod_{D_x^\times}D_x)$）在 $\Bun_G$ 上层范畴（即 $A(\Sigma)$）上的作用，应当在 $x\in\Sigma(F)$ 上局部常值。并非所有 $\Bun_G$ 上的层都满足这一要求，因此我们应当替换 A 侧，考虑子范畴 $\Shv_N(\Bun_G)$，使这种局部常值性成立。于是 Betti 或 étale 设定下修改后的几何 Langlands 对应陈述是一个等价
\[
  \mathbb L_G:\ \Shv_N(\Bun_G)\to\QC^!(\Loc_{\check G}).
\]

\begin{rem}
在 de Rham 设定中不需要条件 $N$，因为每个 D-模都在较弱的意义下局部常值。
\end{rem}

下面两个 $G=\mathbb G_m$ 的例子从两个角度说明条件 $N$ 的必要性。两个例子都在 Betti 设定下工作。

\begin{exmp}[TQFT]
令 $G=\mathbb G_m$。我们有 $\Bun_G=\Pic$。给定 $n\in\mathbb Z$ 与点 $x\in\Sigma(F)$，有自同构
\[
  m_x^n:\ \Pic\to\Pic,\qquad \mathcal L\mapsto\mathcal L(nx).
\]
把 $\Rep\mathbb G_m^\vee$ 的权 $n$ 表示看作点 $x\in\Sigma(F)$ 处球 Hecke 范畴中的对象。该对象在 $\Shv(\Pic)$ 上的作用等同于 $(m_x^n)^*:\Shv(\Pic)\to\Shv(\Pic)$。空间 $\Pic$ 有 $\mathbb Z$ 个分支，按次数标号；由 Hecke 作用的局部常值条件可知，$F\in A_G(\Sigma)$ 应当在每条分支上局部常值。
\end{exmp}

\begin{exmp}[直接计算]
回忆 $\Pic=\Jac\times(\cdots)$、$\Loc_{\mathbb G_m}=H^1(\Sigma,\mathbb C^\times)\times(\cdots)$，其中 $(\cdots)$ 是某些额外因子。为简单起见，把两侧的额外因子都丢掉。我们有
\[
  \QCoh(H^1(\Sigma,\mathbb C^\times))=O(H^1(\Sigma,\mathbb C^\times))\text{-}\mathrm{Mod}.
\]
Fourier 变换给出同构 $O(H^1(\Sigma,\mathbb C^\times))\cong\mathbb C[H_1(\Sigma,\mathbb Z)]$，右端应理解为群代数。于是得到
\[
  \QCoh(\Loc_{\mathbb G_m})\cong H_1(\Sigma,\mathbb Z)\text{-}\mathrm{Mod}.
\]
另一方面 $H_1(\Sigma,\mathbb Z)=\pi_1(\Jac)$，于是我们确认了等价
\[
  \Loc(\Jac)\cong\QCoh(\Loc_{\mathbb G_m}).
\]
由此可以断定：几何 Langlands 对应迫使 A 侧的层局部常值。这与上一个例子的结论一致。
\end{exmp}

到此我们已经修改了几何 Langlands 对应。要陈述 BZSV 猜想还得再作一处修改，因为周期层不一定属于范畴 $\Shv_N(\Bun_G)$。不过我们有 Beilinson 投影 $(-)_{\mathrm{spec}}:\Shv(\Bun_G)\twoheadrightarrow\Shv_N(\Bun_G)$，于是可以期待下面这一点。

\begin{conj}[整体几何猜想]
周期层 $P_X$ 的投影与 L-层 $L_{\check X}$ 在几何 Langlands 对应下匹配：
\[
  \mathbb L_G:\ \Shv_N(\Bun_G)\to\QC^!(\Loc_{\check G}),\qquad (P_X)_{\mathrm{spec}}\mapsto L_{\check X}.
\]
\end{conj}

\begin{rem}
我们忽略了一些技术问题，例如归一化与对偶对合。该猜想更精确的形式见 BZSV 的猜想 12.1.1。
\end{rem}

\subsection{关于 Beilinson 投影}

为说明 Beilinson 投影对周期层的必要性，我们来看周期层与 L-层的一些例子。

\begin{exmp}
令 $X=\mathrm{pt}$，我们称此情形为 Neumann 边界条件。此时映射 $\pi:\Bun^X_G\to\Bun_G$ 是恒等，故 $P_{\mathrm{pt}}=k_{\Bun_G}$。若令 $\check X=\mathrm{pt}$，同样地 $L_{\mathrm{pt}}=\omega_{\Loc_{\check G}}$。注意 $\omega_{\Loc_{\check G}}$ 有来自 $\Loc_{\check G}$ 的辛结构的平凡化。这个辛结构用 $\mathfrak g^\vee$ 的 Killing 型典范地定义。
\end{exmp}

\begin{exmp}
若取 $G$-空间 $X=G$，我们称此情形为 Dirichlet 边界条件。此时映射 $\pi:\Bun^X_G\to\Bun_G$ 是平凡丛处的包含 $\mathrm{pt}\hookrightarrow\Bun_G$。于是 $P_G$ 是平凡丛处的摩天楼层。同样地，$L_{\check G}$ 是平凡局部系处的摩天楼层。
\end{exmp}

\begin{rem}
若 $G$ 交换，这两个例子经 Fourier 变换互为对偶。
\end{rem}

在例 13.24 中，周期层 $P_G$ 是摩天楼层，因而远非局部常值。这告诉我们周期层不一定住在范畴 $\Shv_N(\Bun_G)$ 里，从而说明了 Beilinson 投影的必要性。

\begin{rem}
L-层一般并不落在范畴 $\IndCoh_{\mathrm{Nilp}}(\Loc_{\check G})$ 中。注意我们用的是「更大」版本的几何 Langlands 对应，允许 B 侧上一切 ind-凝聚层。
\end{rem}

条件 $N$ 与 Beilinson 投影看起来像是整体几何猜想的讨厌特征。不过我们关心的 $\Shv_N$ 中的测量可以通过伴随提升为 $\Shv$ 中的测量。投影 $(-)_{\mathrm{spec}}:\Shv(\Bun_G)\twoheadrightarrow\Shv_N(\Bun_G)$ 与包含 $\Shv_N(\Bun_G)\hookrightarrow\Shv(\Bun_G)$ 之间有伴随。在 Betti 设定中，投影左伴随于包含；在 étale 设定中，投影右伴随于包含。用这个伴随，在 Betti 设定中有
\begin{equation}
  \Hom_{\Shv_N}((P_X)_{\mathrm{spec}},F)=\Hom_{\Shv}(P_X,F)
  \label{eq:13-27}
\end{equation}
其中 $P_X$ 是周期层，$F$ 是 Hecke 特征层。若把 $\Hom$ 空间反过来看，étale 设定也有类似陈述。注意在数论中，周期与自守函数之间的配对才是我们在 A 侧主要关心的测量。因此数论感觉不到范畴 $\Shv$ 与 $\Shv_N$ 的差别。

我们用一个简单例子结束对 Beilinson 投影的讨论。

\begin{exmp}
回忆例 13.24 中 $P_G=i_!k$，其中 $i:\mathrm{pt}\hookrightarrow\Bun_G$ 是平凡丛处的包含。如前所述，它不满足条件 $N$，因而需要投影到范畴 $\Shv_N$。思路是把 $\mathrm{pt}$ 换成一个同伦等价的空间 $C$（即可缩空间），以便沿映射 $j:C\to\Bun_G$ 取 $!$-推前。若选取的替换 $C$ 使 $C\to\Bun_G$ 是纤维化，则层 $j_!k_C$ 局部常值，因而落在 $\Shv_N$ 中。我们令 $(P_G)_{\mathrm{spec}}=j_!k_C$。
\end{exmp}

\subsection{Iwasawa--Tate 情形的合理性检验}

相对朗兰兹对偶最重要的例子之一是 Iwasawa--Tate 情形，即 $G=\mathbb G_m$ 作用在 $X=\mathbb A^1$ 上，$\check G=\mathbb G_m$ 作用在 $\check X=\mathbb A^1$ 上。该情形下猜想 13.21 的证明由 Feng--Wang 的 \emph{Geometric Langlands duality for periods} 定理 6.1.2 给出。我们不重复其证明，而是通过比较两侧的 $\Hom$ 空间作一个简单的合理性检验。本小节只在 Betti 设定下工作。这一检验更详细的叙述见 BZSV 第 12.2.3 节。

设 $E$ 是 $\Sigma$ 上非平凡的 $\mathbb G_m$-局部系，它确定 $\QC^!(\Pic)$ 中的摩天楼层 $C_E$。在几何 Langlands 对应下，它对应 $\Pic$ 上的一个秩一 $\mathbb C$-线性局部系 $F_E$。局部系 $F_E$ 有如下性质：对每个 $r\ge0$，它在映射
\[
  \phi_r:\ \Sigma_r\to\Pic,\qquad (P_1,\cdots,P_r)\mapsto\mathcal O(P_1+\cdots+P_r)
\]
下的拉回是 $E^{\boxtimes r}$。这里 $E$ 也用来记与它相伴的秩一 $\mathbb C$-线性局部系。于是猜想 13.21 预言
\[
  \Hom_{\Shv_N}((P_{\mathbb A^1})_{\mathrm{spec}},F_E)\cong\Hom_{\QC^!(\Loc_{\mathbb G_m})}(L_{\mathbb A^1},C_E).
\]
由伴随 (13.27)，它等价于
\begin{equation}
  \Hom_{\Shv}(P_{\mathbb A^1},F_E)\cong\Hom_{\QC^!(\Pic)}(L_{\mathbb A^1},C_E).
  \label{eq:13-29}
\end{equation}
下面给出此式的证明。

由伴随，(13.29) 的右端是 $L_{\mathbb A^1}$ 在 $E$ 处纤维的对偶。如 13.5 小节所述，$L_{\mathbb A^1}$ 在 $E$ 处的纤维是 $R\Gamma(\check\pi^{-1}E,\omega)$，其中 $\check\pi^{-1}E$ 是 $\pi_1(\Sigma)$ 作用在 $\mathbb A^1$ 上的导出不动点轨迹。由注 13.12，对偶化层的导出截面与 $\check\pi^{-1}E$ 的函数代数等同。合起来得到
\begin{equation}
  \Hom_{\QC^!(\Pic)}(L_{\mathbb A^1},C_E)\cong O(\check\pi^{-1}E)^*.
  \label{eq:13-30}
\end{equation}
由于我们假设 $E$ 非平凡，导出栈 $\check\pi^{-1}E$ 只有一个经典点，即零截面。于是它的函数代数由切复形决定。切复形是 $H^1(\Sigma,E)[-1]$，故
\begin{equation}
  O(\check\pi^{-1}E)=\Sym(H^1(\Sigma,E)^*[1]).
  \label{eq:13-31}
\end{equation}
结合 (13.30) 与 (13.31) 得到
\begin{equation}
  \Hom_{\QC^!(\Loc_{\mathbb G_m})}(L_{\mathbb A^1},C_E)\cong\Sym(H^1(\Sigma,E)[-1]).
  \label{eq:13-32}
\end{equation}

现在计算 (13.29) 的左端。映射 $\check\pi:\Bun^{\mathbb A^1}_{\mathbb G_m}\to\Pic$ 有一个截面 $i:\Pic\to\Bun^{\mathbb A^1}_{\mathbb G_m}$，相当于选取相伴 $\mathbb A^1$-丛的零截面。截面 $i$ 是闭嵌入，于是可以考虑补开浸入
\[
  j:\ \Bun^{\mathbb A^1,0}_{\mathbb G_m}\to\Bun^{\mathbb A^1}_{\mathbb G_m}.
\]
注意开子栈 $\Bun^{\mathbb A^1,0}_{\mathbb G_m}$ 分类的是一对 $(L,s)$，其中 $L$ 是 $\Sigma$ 上的线丛、$s$ 是它的非平凡截面。按重数取截面的零点，这个栈等同于有效除子的栈 $\Sym\Sigma$。于是我们有正合三角
\begin{equation}
  C_{\Pic}\to P_{\mathbb A^1}\to j_!C_{\Sym\Sigma}.
  \label{eq:13-33}
\end{equation}
我们断言 $\Hom_{\Shv}(C_{\Pic},F_E)=0$。

回忆 $\Pic$ 由按 $\mathbb Z$ 编号的多份 $\Jac\times B\mathbb G_m$ 组成（例 13.5），而 Hecke 算子置换这些拷贝（例 13.19）。由于局部系 $F_E$ 是非平凡的 Hecke 特征层，上述事实告诉我们 $F_E$ 在每份 $\Jac\times B\mathbb G_m$ 上都非平凡。于是为证明该断言，只需证明：对 $\Jac\times B\mathbb G_m$ 上任何非平凡秩一局部系 $F$，有 $\Hom_{\Jac\times B\mathbb G_m}(C_{\Jac\times B\mathbb G_m},F)=0$。这等价于 $F$ 的单值表示的群上同调的一个消没陈述。由于 $\Jac$ 是环面，该断言由下述事实推出：对有限秩自由交换群 $\Lambda$ 及其非平凡特征 $k_\chi$，群上同调群 $H^*(\Lambda,k_\chi)$ 消没。

把该断言与正合三角 (13.33) 结合，得到
\[
  \Hom_{\Shv}(P_{\mathbb A^1},F_E)\cong\Hom_{\Shv}(j_!C_{\Sym\Sigma},F_E).
\]
用伴随以及 $j$ 是开浸入这一事实，可改写为
\begin{equation}
  \Hom_{\Shv}(P_{\mathbb A^1},F_E)\cong R\Gamma(\Sym\Sigma,j^*F_E).
  \label{eq:13-34}
\end{equation}
回忆 $F_E$ 经 $\phi_r:\Sigma_r\to\Pic$ 拉回为 $E^{\boxtimes r}$，故
\[
  R\Gamma(\Sigma_r,\phi_r^*F_E)=R\Gamma(\Sigma,E)^{\otimes r}=(H^1(\Sigma,E)[-1])^{\otimes r}.
\]
由此可以期待
\[
  R\Gamma(\Sym\Sigma,j^*F_E)\cong\Sym(H^1(\Sigma,E)[-1]).
\]
这与 (13.34) 结合即给出
\begin{equation}
  \Hom_{\Shv}(P_{\mathbb A^1},F_E)\cong\Sym(H^1(\Sigma,E)[-1]).
  \label{eq:13-35}
\end{equation}
比较 (13.32) 与 (13.35)，我们通过了合理性检验 (13.29)。
